\head{Dot续译整合核查：预备笔记第7--11节}
\source{\texttt{EGAVaiello.pdf}第19--43页；十八批Dot实存译稿与本版补译第11节}

\subsection*{（7.4）的反例}
\noindent\textbf{（7.4）按字面的反例。}
取 $k=\mathbf F_p(t)$ 及 $L=k(u)$，其中 $u^p=t$。扩张 $L/k$ 是次数 $p$ 的非平凡纯不可分扩张。$X=\Spec L\to\Spec k$ 有限、平坦、固有，且 $\dim X=0$；其不可分点集也是维数0，所以符合底本的 $\leq n$。

若存在有限平展满射 $S'\to\Spec k$，则 $S'=\Spec A$，其中 $A$ 为非零有限可分 $k$-代数，即有限个有限可分域扩张 $K_i/k$ 的直积。$X_{S'}\to S'$ 的截面会诱导一个 $k$-代数同态 $L\to A$，再投影到任一 $K_i$，得到 $k$-嵌入 $L\hookrightarrow K_i$。这是不可能的，因为有限可分扩张的子扩张可分，而 $L/k$ 非平凡纯不可分。由此，所印条件不能保证结论。此处译校不自行把 $\leq n$ 改成 $<n$，也不声称仅改变这一符号就能修复全部证明。



\subsection*{余法丛论证的补充核验}
正合列的源靶、对偶Euler列、法模与余法模的区别、射影化约定、$(1,0)$ 与 $(1,1)$ 扭转的区别。为核实“机械性的验证”，可在局部平凡化下写 $\phi=\sum_i a_i u_i$。相对于 $P^\vee$ 作微分时，$a_i$ 为常量，所以 $d\phi=\sum_i a_i\,du_i$；在 $Y$ 上关系 $\sum_i a_i u_i=0$ 使 $L\to E_Z$ 经由Euler列的核分解，而复合到 $\Omega^1_{X/S}$ 正给出上式。这是译校补充验证，不冒称原稿展开了这一步。




\subsection*{二次覆盖反例的独立核验}
\subsubsection*{次数二的反例逐步核验}
令 $s=t^2$。在 $\mathbb C$ 上，$\mathbb G_m\to\mathbb G_m$ 的微分为 $2t\,dt$，处处非零，且该映射有限平展、次数二。映射 $s\mapsto[1:s:s^2]$ 是非退化光滑圆锥曲线的一个开子集的闭嵌入到相应开集。故复合 $f$ 是非分歧态射；$X$ 光滑且整，符合现存背景。

圆锥曲线在 $[1:s:s^2]$ 的切线方程为
\[
s^2X_0-2sX_1+X_2=0.
\]
拉回到参数 $u\in X$，其方程成为
\[
t^4-2t^2u^2+u^4=(u^2-t^2)^2.
\]
在 $u=t$ 和 $u=-t$ 处，它均为非退化二重零点，因为局部二次项系数 $4t^2$ 非零。这也直接说明一个切超平面的非超奇异接触点可以有两个。余法模秩为一，故 ${Y^{\mathrm{sing}}}\simeq X$；参数化的像闭包满足 $b^2=4ac$，为光滑对偶圆锥曲线。映射 $g$ 的函数域扩张为 $\mathbb C(t)/\mathbb C(t^2)$，不是次数一。因此反例不依赖对残缺切空间句子的任何猜测。


\subsection*{D17-E01：（7.4）原断言与证明的完整对照}
\noindent\textbf{以下是底本按字面的旧译，不作为有效命题采用。}
\noindent\textbf{原命题（7.4）。}
设 $g:X\to S$ 为平坦固有态射，$s\in S$。令 $n=\dim X_s$，并假设 $X_s$ 中 $X_s$ 在 $k(s)$ 上不可分的点所成集合的维数 $\leq n$，例如 $X_s$ 可分。那么，存在 $s$ 的开邻域 $U$，以及平展、有限且满射的态射
\[
 S'\longrightarrow U,
\]
使得 $X\times_S S'$ 在 $S'$ 上有一个截面。如果 $k(s)$ 是无限域，则可以取 $S'$ 为 $X|_U$ 的一个闭子概形。\ednote{不清楚，请问 A.G.。〔底本脚注27。〕}
\begin{proof}
先假设 $k(s)$ 无限。对 $n$ 归纳，$n=0$ 的情形是平凡的。事实上，这时存在 $s$ 的一个开邻域 $U$，使 $X|_U$ 本身在 $U$ 上平展、有限且满射；通过直接引用有关结果即可看出这一点。如果 $n>0$，把（7.3）应用于“可分”这一性质，就能用一个具有同样性质的“超平面截面”$Y$ 来替代 $X$，只是把 $n$ 换成 $n-1$。如果不假设 $k(s)$ 无限，则先作一个平展基变换，论证仍然适用。（可以进行下去。）
\transnote{底本写的是 $\leq n$，不是 $<n$。这一条件不能支撑其写出的 $n=0$ 起始步。例如 $X=\Spec(k[\varepsilon]/(\varepsilon^2))$、$S=\Spec k$，其中 $k$ 为无限完美域，则 $X\to S$ 平坦且有限，故固有；$n=0$，不可分点集的维数至多为0。但这个非约化的 $X$ 不在 $k$ 上平展。因此“$X|_U$ 本身平展”的推断不成立。此例只反驳所写证明的这一推断，并不反驳这个例子中的截面存在结论。另取无限的不完美域 $k=\mathbf F_p(t)$ 与 $L=k(t^{1/p})$，令 $X=\Spec L$；同样满足所写的维数上界，而不存在非空有限平展 $k$-概形 $S'$ 使 $X_{S'}\to S'$ 有截面，详见此处译校后附核查说明。故所印条件下的命题本身也不能直接使用。}
\end{proof}



\subsection*{D17-E02：（8.4）及（8.4.1）旧稿对照}
\noindent\textbf{以下保留旧译所见的错误符号与残句，仅供核对；现正文已订正确定的维数和包含关系。}
\noindent\textbf{原命题（8.4）。}
设 $T$ 是 $X$ 的闭子集，并假设 $\codim(T,X)\geq k$。那么，对每个 $\xi\in P^\vee$，都有 $\codim(T_\xi,Y_\xi)\geq k-1$。令 $Z$ 为满足
\[
\codim(T_\xi,Y_\xi)=k-1
\quad\text{（即 }\codim(T_\xi,Y_\xi)<k\text{）}
\]
的 $\xi\in P^\vee$ 所成的集合，则 $Z$ 是 $P^\vee$ 的可构造、无处稠密〔法文：稀疏〕子集，即 $\overline Z$ 在 $P^\vee$ 中的余维至少为一。

它的余维至少为二的充要条件是：对 $X$ 的每个余维等于 $k$ 且满足 $\dim f(T_i)=0$ 的不可约分支 $T_i$，都存在 $X$ 的一个不可约分支 $X_j$，使
\[
\codim(T_i,X_j)=k,\qquad \dim f(X_j)-0
\]
〔或者“闭包”被划去了？〕；即，若 $f$ 拟有限且 $k>0$，则〔原文有不明确的 que (it???)〕$T$ 没有满足 $\dim_xX-k$ 的孤立点。第一个断言立即由下面的引理（8.4.1）(a) 得出；这个引理是对第5节的一项事后补充。
\transnote{此处 $T_i$ 确写为 $X$ 的分支，而不是 $T$ 的分支；当 $k>0$ 时，这与其在 $X$ 中余维等于 $k$ 的要求冲突。减号写法按底本保留，不据猜测填成等号。}

\par\noindent\textbf{原引理（8.4.1）。}
设 $X$ 为局部诺特预概形，$L$ 为 $X$ 上的可逆模，$\phi$ 为 $L$ 的一个截面，$Y=V(\phi)$，$T$ 为 $X$ 的闭子集。假设
\[
\codim(Y,X)\geq k\text{。}
\]
则：
\begin{enumerate}
\item[(a)] $\codim(T\cap Y,Y)\geq k-1$。
\item[(b)] 要有
\[
\codim(T\cap Y,Y)=k-1,
\quad\text{即 }\codim(T\cap Y,Y)<k,
\]
其充要条件是：存在 $T$ 的一个包含于 $Y$ 的不可约分支 $T_i$，使
\[\codim(T_i,X)=k,\]
并且对 $X$ 的每个包含 $T_i$ 且满足
\[
\dim\mathcal O_{X_j,T_i}=\dim\mathcal O_{X,T_i}\quad(=k)
\]
的不可约分支 $X_j$，都有 $X_j\not\subset Y$。
\end{enumerate}
根据 $0_{\mathrm{IV}}$、第IV章中关于维数的一般事实，这个引理的验证是直接的。
\transnote{假设中的 $Y$ 按底本保留。此假设下(b)有明确反例：取 $X=\mathbb A^1$、$L=\mathcal O_X$、$\phi=t$、$Y=\{0\}$、$T=X$、$k=1$。假设成立，且 $\codim(T\cap Y,Y)=0=k-1$；但 $T$ 唯一的不可约分支是 $X$，并不包含于 $Y$，所以所声称的必要条件失败。疑应以 $\codim(T,X)\geq k$ 为假设，但此处译校不据此替换或声称已经证明修正版。}

在（8.4）的假设下，由（8.4.1）(b) 可以看出哪些超平面 $H_\xi$ 是例外的。排除如下 $\xi\in P^\vee$ 所成的集合 $Z_0$：$T$ 或 $X$ 有一个不可约分支 $R$，满足 $\dim f(R)>0$，且 $R_\xi$ 的“维数过大”。这是一个余维为二的集合，在以下论述中不予计入。例外的 $H_\xi$ 是那些使得存在一个 $T_i$ 满足\ednote{大概是 $H_\xi$。〔底本脚注32。〕}
\[
\codim(T_i,X)=k,\qquad\dim f(T)=0,\qquad f(T)\ \mathrm{CH}
\]

并且，对 $X$ 的每个满足 $X_j\supset T_i$ 及 $\codim(T_i,X_j)=k$ 的不可约分支 $X_j$，都有 $f(X_j)\not\subset H_\xi$ 的超平面。

给定满足 $\codim(T_i,X)=k$ 的 $T_i$，若存在 $X_j$，使 $\codim(T_i,X_j)=?$〔字迹不清，请询问 Grothendieck〕且 $\dim f(X_j)=0$，则有 $f(X_j)=f(T_i)\not\subset H_\xi$，所以 $\xi$ 相对于 $T_i$ 不是例外的。反之，如果对每个包含 $T_i$ 且满足 $\codim(T_i,X_j)=k$ 的 $X_j$ 都有 $f(X_j)>0$，则对 $\xi\in P^\vee-Z_0$，$\xi$ 相对于 $T_i$ 为例外的充要条件是 $f(T_i)\not\subset H_\xi$；这样的 $\xi$ 所成的集合，是 $P^\vee$ 中一个超平面在 $P-Z_0$ 上的迹。这就证明了（8.4），并且证明了更精确的结果：例外集是一个余维至少为二的集合与若干超平面的并；这些超平面可由以上证明明显地确定。
\transnote{保留了 $f(T)\ \mathrm{CH}$、$f(X_j)>0$、$P-Z_0$ 以及未辨条件，没有补齐。末尾的“不包含”与“超平面”冲突：固定有理点不在 $H_\xi$ 中的参数集合是开集；使超平面包含该点的参数集合才是对偶超平面。}

（恐怕这段写法很松散，或者说很潦草〔英译者措辞说明〕，因为我一直在几何意义下推理，却没有说出来，即取值于代数闭域的点。当然，（8.4）中所陈述的条件确实是几何的，所以可以假设 $k$ 代数闭，并对 $k$-有理点进行论证。）利用（8.4）、（5.7.4）以及（8.2）的末尾，得到：


\subsection*{D17-E03：（10.1）例外集与闭包}
\noindent\textbf{旧稿按底本保留的整段：}
设 $Z$ 是 $P^v$ 中相对于一个可构造性质 $\mathcal P$ 的 $\mathcal P$-例外集，使得 $Z$ 是 $P^v$ 的可构造子集。设 $S=\operatorname{Spec}(k)$。我们将看到（参见第12节，那里将补述一些无疑本应在前几节出现的内容），要有
\[
\operatorname{codim}(Z,P^v)\ge2,
\]
〔以下为原稿断言，其中单条直线的充分条件待核。〕必要且充分的是：$P^v$ 中“每一条充分一般的直线” $L$ 都不与 $Z$ 相交，或者也可以说（甚至）不与 $Z$ 相交；而只要存在一条（单独一条）不与 $Z$ 相交的 $L$ 就足够了〔底本编辑批注：“应为 $Z$；A.G.的错误，P.B.”〕。\footnote{原脚注55：“请询问A.G.”}\transnote{这一整句有未解决的排印或编辑问题：底本在“or even”前后及“should be”前后均印作 $Z$，未见可辨的上划线，而下一句确实印为 $\bar z$。这里按实际可辨文字保留重复，不把它擅自恢复成关于 $\overline Z$ 的定理。对一般可构造集，“存在一条直线不交 $Z$”不足以推出余维至少2，例如在射影平面中从一条直线删去一点所得的集合；经过该点的另一条直线可避开该集合。故该处不能被当作已核实的充分条件。底本此处及后文数次写 $L\subset P$，但开头及10.2明定参数直线在对偶空间 $P^v$ 中；正文在这些明确的参数空间处统一为 $P^v$。}
若 $k$ 无限，则必要且充分的是存在一条不与 $\overline Z$ 相交的直线 $L\subset P^v$。\transnote{底本印作小写 $\bar z$；上下文没有引入另一个 $z$，此处按同一例外集统一为 $\overline Z$。这是原稿对上述余维条件的表述，未在此处译校中独立证明；“充分一般”及直线的定义域约定仍须结合第12节核查。}

\subsection*{D17-E02：修正（8.4.1）后的独立补证}
底本第23页将假设写成 $\codim(Y,X)\geq k$；它已有前述 $X=\mathbb A^1$、$Y=\{0\}$、$T=X$ 的反例。正文改用 $\codim(T,X)\geq k$，以下验证修正后的两项，而不是仅据上下文猜换字母。

对 $z\in T\cap Y$，取局部环 $R=\mathcal O_{X,z}$，并在平凡化 $L$ 后以 $f\in R$ 表示 $\phi$。因 $z\in T$，有 $\dim R\geq k$。局部诺特环中一个方程使维数至多降低一，故
\[
\dim R/(f)\geq\dim R-1\geq k-1.
\]
在 $T\cap Y$ 的泛点取下确界，即得(a)。当下确界恰为 $k-1$ 时，有某个这样的泛点 $z$ 满足 $\dim R=k$。如果 $z$ 严格特化于包含它的某个 $T_i$ 的泛点，那么在该泛点下方长至少 $k$ 的素理想链后，再接上这个严格特化，就给出 $\dim R\geq k+1$，矛盾。所以 $z$ 正是 $T$ 的某个分支 $T_i$ 的泛点，且 $T_i\subset Y$、$\codim(T_i,X)=k$。

在这个维数 $k$ 的局部环 $R$ 中，$\dim R/(f)=k-1$ 当且仅当 $f$ 不属于任何满足 $\dim R/\mathfrak p=k$ 的极小素理想 $\mathfrak p$。若属于其中一个，商环维数仍至少为 $k$；若均不属于，则每个最高维不可约分支的非零主理想商维数至多为 $k-1$，其他分支本来维数就不超过 $k-1$，再与上述下界合用即得等号。这些极小素理想正对应正文(b)中满足局部维数条件的 $X_j$；$f\notin\mathfrak p$ 正是 $X_j\not\subset Y$。反向按同样的局部论证成立，故(b)得证。

同时纠正（8.4）证明末尾的一个强于论证的断言。例外集一般只在除去余维至少2的集合后，等于若干对偶超平面的迹；不能不加限定地说它是若干整条超平面与低维集合的并。取 $X\subset\mathbf P^2$ 为一条直线、$T=\{x\}$、$k=1$。经过 $x$ 而不等于 $X$ 的直线 $H$ 给出 $Y_H=T$，所以 $\codim(T_H,Y_H)=0$；但 $H=X$ 时 $Y_H=X$，余维为1。因此例外集是一条对偶直线删去一点。这一实例同时核验了正文用“在开集上的迹”替代原精确并集句的必要性。

\subsection*{D17-E03：（10.1）对闭包的修正}
令 $Z=H\setminus\{x\}\subset\mathbf P^2$，其中 $H$ 是直线、$x\in H$。任取经过 $x$ 而不等于 $H$ 的直线 $L$，就有 $L\cap Z=\varnothing$，但 $\overline Z=H$ 的余维为1。这反驳“存在单条不交可构造集 $Z$ 的直线，便有余维至少2”。对于闭集 $C$，任一余维1分支是射影超曲面，必与每条几何直线相交；若 $\codim C\geq2$，与 $C$ 相交的直线在直线 Grassmann 簇中构成真闭集（可用点—直线关联空间及维数界计算），故一般直线避开 $C$。在无限域上，Grassmann 簇的标准仿射图中非空开集有有理点，于是可选 $k$-有理直线。这给出正文所列的精确区分。

\subsection*{D17-E04--E06：§11 补译与明确订正}
Dot 工程的（10.4）批在进度中指定下一处为底本第41页§11，而后一个可恢复批次从第43页§12开始；54个旧批次内没有§11主源码。此版对照所上传 \texttt{EGAVaiello.pdf} 第41--43页逐页图像，补译§11全部段落及原脚注57--60。其起点为 Grassmanians 标题，终点为 $S$ 不连通时秩可变的附记，不跨入§12。

Plücker 映射必须把 $E_{S'}\twoheadrightarrow G$ 对应到 $\bigwedge^nE_{S'}\twoheadrightarrow\bigwedge^nG$，底本说明句漏掉两处外幂；正文补齐。一般拟凝聚模的 Grassmann 概形可表且分离，并不因是射影化的闭子概形就自动拟紧：取 $E=\bigoplus_{i\geq0}ke_i$、$n=1$，在 $\operatorname{Proj}k[T_0,T_1,\ldots]$ 的标准开覆盖中，任何有限个 $D_+(T_i)$ 都遗漏一个只使另一坐标非零的点。有限型 $E$ 则局部由有限个生成元给出有限标准开覆盖，正文因此将拟紧性移入有限型条件。

对一个坐标分解 $E=E'\oplus E''$，要求 $E'\to G$ 为同构的商，经该同构可唯一写成 $(1,a):E'\oplus E''\to E'$，其中 $a:E''\to E'$。因此开图的参数为 $\mathscr H\!om(E'',E')$ 的截面，在 $\mathbf V(M)=\operatorname{Spec}\operatorname{Sym}M$ 的约定下，代表概形正是 $\mathbf V(\mathscr H\!om(E',E''))$。若仅给未选分裂的短正合列，可局部如此描述，但改变分裂使参数发生平移；开图是相应向量丛的挠子。原文先从基选坐标分解，故正文保留这项已给的分裂，避免把后句读成任意正合列的典范向量丛同构。

有限表示的下降通过一个有限表示矩阵实现：取其系数生成的有限型子环 $A_0\subset A$，同一矩阵定义有限表示 $A_0$-模 $E_0$，且 $E_0\otimes_{A_0}A=E$。$A_0$ 诺特，所以对应的有限型射影 Grassmann 概形为有限表示，再经基变换得到结论。底本仅写“有限型模”，正文补明这里所用的是有限表示矩阵；不将任意无限关系的有限型模宣称为可由有限系数下降。

其余订正包括 $E'$ 的首项撇号、分解标记 $s/S$、重复句、双对偶线性子空间的维数排式、以及 $L^0$ 的点误写为 $L$ 的点。正文和脚注分别给出正确写法及原写法。秩超过 $E$ 的秩时 Grassmann 概形为空；不固定秩的函子仍只取局部有限秩商。这些边界说明不表示补出了作者尚未陈述的“次数一”内在刻画。
